\begin{resumo}
Este esqueleto didático apresenta o projeto de pesquisa sobre construção, visibilidade e mobilização estatal de indicadores de violência letal no Brasil, considerando os anos-base de 2016 a 2024. A pesquisa pergunta de que forma escolhas de construção e apresentação desses indicadores se relacionam com sua visibilidade institucional e midiática e com sua mobilização em processos federais de atenção e priorização estatal. O objetivo consiste em analisar essas relações por meio de um desenho documental, longitudinal e de abordagem mista. A fundamentação articula comensuração, comparação entre arenas de atenção e mobilização política de indicadores. O percurso distingue registros administrativos, produtos de consolidação e análise, comunicação institucional, seleção jornalística e instrumentos do Legislativo e do Executivo federal. A proposta preserva uma camada original e uma camada analítica dos dados e não trata divergências entre fontes como erros automáticos. A cobertura jornalística será observada em G1, Folha de S.Paulo e O Estado de S.Paulo, sem pretensão de representar estatisticamente toda a imprensa. O produto técnico esperado é um protocolo de rastreabilidade, comparabilidade e comunicação responsável dos indicadores. O documento demonstra o uso do template LaTeX do IDP, com citações e referências cruzadas, sem apresentar resultados empíricos ou inferir causalidade a partir de correspondências documentais.
\end{resumo}
